\section{Introduction}
\label{sec:intro}

Multi-teacher on-policy distillation (MOPD) transfers domain specialists
into one student by scoring student-generated responses with their
assigned teachers \citep{ma2026mopdmultiteacheronpolicydistillation}.
Its intended outcome is balanced capability integration. Open-MOPD
identifies capability imbalance in a controlled routed setting and
studies token-share weighting, changing distillation gaps, and stale
student-dependent rewards \citep{gao2026openmopddiagnosingfixingcapability}.
These findings motivate examining both the loss weights given to each
domain and the parameter updates induced by its teacher.

We organize the paper around three studies. The first examines loss
reduction as an explanation for imbalance. A global token mean gives
long-response domains greater weight. Averaging tokens separately within
each domain changes those domain weights, but still weights longer
responses more heavily inside the domain. We derive these effects and
test their relationship to multi-task capability. This analysis does
not require a new sampling or weighting algorithm.

The second study is the empirical center of the paper. Prior work
already reports sparse parameter changes in RL
\citep{mukherjee2025subnetworks} and in OPD
\citep{yu2026denseupdates,shen2026opdgeometry}, including support reuse
when varying the teacher. We first verify update sparsity in our
single-teacher and multi-teacher settings. We then ask whether different
teachers acting on the \emph{same MOPD student} update overlapping
parameter subsets. Finally, we test whether teacher pairs with larger
distribution differences induce more different update supports.
Teacher--student distance is measured as a related factor that may help
explain transfer and update differences. Reproducing existing sparsity
results is a foundation for this study, not a claim of first discovery.

The third study is a focused comparison of loss choices. We compare
sampled-token policy-gradient, top-$k$, and full-vocabulary losses in
both single-teacher and multi-teacher OPD, keeping the compared
update measurements consistent. If these choices
show no material sparsity difference in the tested regime, that is a
reportable observation and the study stops at that comparison. If a
stable difference appears, additional controls can examine its source.
We do not make estimator-noise analysis the main contribution or assume
that sampled-token supervision must produce sparser parameter updates.

The intended contributions are therefore:
\begin{enumerate}
\item An empirical and algebraic explanation of how loss reduction
and length weighting affect capability balance across tasks.
\item A study of sparse OPD updates, the overlap of teacher-conditioned
supports within MOPD, and their relationship to teacher--teacher and
teacher--student distribution differences.
\item A comparison of parameter-update sparsity across sampled-token, top-$k$, and full-vocabulary supervision in
single-teacher and multi-teacher distillation.
\end{enumerate}

The core questions permit positive, negative, and null relationships.
Support overlap does not by itself imply harmful interference, and a
sparsity difference does not by itself justify choosing a loss. Capability
and signed cross-task effects provide context for these measurements.
\missing{Replace the prospective empirical contributions with observed
findings and their measured scope once the experiments are complete.}
